% ----------------------
\subsection{Section 83 (30 April 1832)}
% ----------------------
	\label{DC83}
	
	% -----
	\subsubsection{Historical Background}
	% -----
		\label{DC83-HB}
		
		% --
		\paragraph{JSP}
		% --
		
			``According to a later JS history, JS `sat in council with the brethren' on 30 April 1832 in Independence, Jackson County, Missouri, and dictated a revelation clarifying the rights of women and children who had lost their husbands or fathers. Although there were civil laws outlining specific property rights for women upon the death of their husbands, it was not clear what would happen if a husband had consecrated property to the church. The 30 April revelation helped clarify the church's position in such instances.

			``An earlier revelation on the laws of the church outlined principles of consecration, stating that consecration was a means for church members to take care of the poor among them. That revelation instructed members to consecrate their properties to the church and bishops to convey back stewardships that were sufficient for their needs and those of their families. The `residue' of the consecrated property was then kept in the `store house to administer to the poor and needy.' In accordance with this revelation, a few church members in Ohio consecrated properties to the church for a brief time in 1831. After Missouri was designated as the location of Zion in July 1831, the practice of consecration was instituted there. As more members consecrated property, questions inevitably arose. For example, according to Missouri statutes, a woman had no claim on her personal property when she married, giving up that right to her husband. If her husband died, a woman had a dower right, or a right to a third part of her husband's real estate. She was also `entitled absolutely to a share' in her husband's `other personal estate' that was `equal to the share of a child of such deceased husband, after the payment of debts.' The doctrine of consecration, as established in the `Laws of the Church of Christ' and subsequent revelations, did not address what claims a widow had on property her husband had consecrated to the church or what would happen to children who lost their fathers.

			JS apparently became concerned about such questions while in Missouri in April 1832. Perhaps one reason was his short trip from 28--29 April to visit the Saints from Colesville, New York, who had settled about twelve miles west of Independence in Kaw Township, Missouri. JS had good friends in this settlement, including the Knight family. He later reported that he `received a welcome only known by brethren and sisters united as one in the same faith.' Among those living in Kaw Township were at least two widows: Phebe Crosby Peck, who had four children, and Anna Slade Rogers, who had a daughter. These women's husbands died in 1829 before the revelation on the `Laws of the Church of Christ' was dictated, but JS's association with them may have prompted him to wonder about a widow's claim to consecrated property, which may in turn have led to this 30 April revelation.

			Sidney Rigdon may have written this revelation as JS dictated it in late April 1832. The apparent misdating of `May 31\supers{st}' indicates that the version featured here, in Rigdon's handwriting, is probably not the original but is a copy made for Newel K. Whitney. Whitney corrected this error in an endorsement that labeled and characterized the revelation: `as to Women \& children; Inheretance at Zion 30 ap\supers{l.} 1832.' Sometime after April 1832, John Whitmer copied the revelation into Revelation Book 1, giving it the date of 30 April---a date that was perpetuated in a later JS history. It was also published in the January 1833 issue of \textit{The Evening and the Morning Star} under the title `Items in Addition to the Laws for the Government of the Church of Christ, Given April, 1832.'''
			
	% -----
	\subsubsection{Versions}
	% -----
		\label{DC83-V}
		
		See the \href{https://drive.google.com/file/d/1goAvNz8jS4R8slYfMG_db55Ty2z_vmgK/view?usp=sharing}{sections comparison} document.
	
		\begin{compactdesc}
			\item[JSP] \hyperref[EMS-1832-JS-30-APR-1832]{30 April 1832}
			\item[EMS] 1:8 (January 1833), 62
			\item[BC] ---
			\item[DC 1835] Section 88, 222--223
			\item[TS] 5:16 (2 September 1844), 625
			\item[DC 1844] Section 89, 344--345
			\item[MS] 14:11 (8 May 1852), 163
		\end{compactdesc}

	% -----
	\subsubsection{Section 83:1} % *DC83-1 
	% -----
		\label{DC83-1}

		VERILY, thus saith the Lord, in addition to the laws of the church concerning women and children, those who belong to the church, who have lost their husbands or fathers:

		\begin{framed}
		\scriptsize
			\begin{compactdesc}
				\item[JSP] Verily thus saith the Lord in addition to the laws of the church concerning women and children who belong to the church who have lost their husbands or fathers
				% \item[BC] 
				% \item[]
			\end{compactdesc}
		\end{framed}

	% -----
	\subsubsection{Section 83:2} % *DC83-2 
	% -----
		\label{DC83-2}

		Women have claim on their husbands for their \hyperref[MAINTENANCE]{maintenance}, until their husbands are taken; and if they are not found transgressors they shall have fellowship in the church.

		\begin{framed}
		\scriptsize
			\begin{compactdesc}
				\item[JSP] \sout{<a>} \sout{women} \sout{<woman>} \sout{have} \sout{<has>} <have> claim on their husbands untill <t>he<y> \sout{is} <are> taken
				% \item[BC] 
				% \item[]
			\end{compactdesc}
		\end{framed}

	% -----
	\subsubsection{Section 83:3} % *DC83-3 
	% -----
		\label{DC83-3}

		And if they are not faithful they shall not have fellowship in the church; yet they may remain upon their inheritances according to the laws of the land.

		\begin{framed}
		\scriptsize
			\begin{compactdesc}
				\item[JSP] and if they are not found transgressors they remain upon their inheritinces
				% \item[BC] 
				% \item[]
			\end{compactdesc}
		\end{framed}

	% -----
	\subsubsection{Section 83:4} % *DC83-4 
	% -----
		\label{DC83-4}

		All children have claim upon their parents for their \hyperref[MAINTENANCE]{maintenance} until they are of age.%
			\footnote{%
				% \label{}%
				From the Family Proclamation: ``Parents have a sacred duty to rear their children in love and righteousness, to provide for their physical and spiritual needs, and to teach them to love and serve one another, observe the commandments of God, and be law-abiding citizens wherever they live. Husbands and wives---mothers and fathers---will be held accountable before God for the discharge of these obligations.''
				}

		\begin{framed}
		\scriptsize
			\begin{compactdesc}
				\item[JSP] all children have claim upon their parents untill they are of age
				% \item[BC] 
				% \item[]
			\end{compactdesc}
		\end{framed}

	% -----
	\subsubsection{Section 83:5} % *DC83-5 
	% -----
		\label{DC83-5}

		And after that, they have claim upon the church, or in other words upon the Lord's storehouse, if their parents have not wherewith to give them inheritances.

		\begin{framed}
		\scriptsize
			\begin{compactdesc}
				\item[JSP] and after that they have claim upon the church or in other words the Lords storehouse for inheritences [\textit{1/4 page blank}]
				% \item[BC] 
				% \item[]
			\end{compactdesc}
		\end{framed}

	% -----
	\subsubsection{Section 83:6} % *DC83-6 
	% -----
		\label{DC83-6}

		And the storehouse shall be kept by the consecrations of the church; and widows and orphans shall be provided for, as also the poor.  Amen.

		\begin{framed}
		\scriptsize
			\begin{compactdesc}
				\item[JSP] ---
				% \item[BC] 
				% \item[]
			\end{compactdesc}
		\end{framed}